\documentclass[a4paper]{article}
\usepackage{graphicx}
\usepackage{amsmath,amssymb}
\usepackage{hyperref}
\usepackage{minted}
\usepackage{multirow}
\usepackage{float}

\title{}
\date{}

\begin{document}
    \maketitle

   
    \section{Surprise in Active Inference: Causes and Modulators }

    ## Executive Summary
    Surprise (self-information) in active inference is the negative log-probability of an observation under the agent’s generative model.  Mathematically, if an agent predicts an outcome with probability \(P(o)\), the surprise is \(I(o) = -\ln P(o)\)【10†L115-L117】.  Because an agent cannot compute this directly (it requires marginalizing over hidden states), it minimizes **variational free energy** \(F\), which satisfies \(F \ge -\ln P(o)\)【10†L43-L49】【27†L470-L478】.  Thus factors that *decrease* \(P(o)\) (making observations improbable) raise surprise and free energy.  Key factors include **model mismatch** (wrong priors or likelihoods), **unexpected observations** (data outside the predicted distribution), and **low prior probability** on the actual causes.  High **precision** (inverse variance) amplifies prediction errors: a small discrepancy is highly surprising if predictions are precise【42†L406-L414】【51†L318-L324】.  In contrast, high sensory noise (low precision) broadens predicted distributions and can *dampen* surprise for a given deviation, although very broad (high-noise) priors raise baseline uncertainty.  In hierarchical models, each level sends predictions downward; mismatches produce **prediction errors** that are weighted by precision (attention).  High-level *expected precision* (attention) further amplifies bottom-up errors, increasing surprise when mismatches occur【42†L406-L414】【51†L318-L324】.
    
    Actions and policy selection aim to minimize **expected surprise** (expected free energy) about future inputs.  Agents choose policies that are expected to produce preferred (high-probability) outcomes or to resolve uncertainty. When actions deliver unpredicted outcomes, unexpected high prediction errors arise, raising surprise and triggering belief updates.  Simple Gaussian examples illustrate these effects: for instance, observing \(y=2\) under a Normal likelihood \(N(\mu,\sigma^2)\) gives surprise \(I(y) = \frac12\ln(2\pi\sigma^2) + \frac{(y-\mu)^2}{2\sigma^2}\).  Decreasing \(\sigma\) (higher precision) or having \(\mu\) far from \(y\) increases \(I(y)\).  Empirical studies support these principles: e.g., increased sensory noise tends to reduce confidence (due to lower precision), but interoceptive arousal (unexpected bodily signals) can modulate this effect【53†L683-L690】.  In schizophrenia, aberrant precision weighting (e.g. overly strong priors or faulty attenuation of sensory signals) leads to excessive prediction errors and ‘spurious’ surprise, contributing to hallucinations and delusions【35†L229-L237】.
    
    **Tables** below compare how various factors (likelihood width, priors, noise, precision) enter the mathematics of surprise and qualitatively alter surprise.  **Figures** illustrate hierarchical predictive coding (flow of predictions, errors, and precision) and the timeline of how surprise arises (from prediction to observation to update).  Together, this analysis shows that *surprise is high whenever observations are unlikely under the agent’s model*, and that modulating beliefs (priors, likelihoods) or observations (via action) can reduce surprise over time.  These principles have direct implications for perception (efficiency vs. novelty), learning (model updates), and psychopathology (e.g. abnormal precision in schizophrenia【35†L229-L237】).
    
    ## Background: Surprise and Free Energy in Active Inference
    In information theory, **surprise** (a.k.a. self-information or surprisal) of an outcome \(o\) is defined as
    \[
    I(o) = -\ln P(o),
    \]
    where \(P(o)\) is the marginal likelihood (evidence) of that observation under the agent’s model【10†L115-L117】.  An event with low probability is highly surprising (large \(I\)), whereas a perfectly predicted event has \(I=0\).  Active inference argues that adaptive agents minimize surprise (on average) to maintain a bounded repertoire of internal states【3†L386-L394】【10†L43-L49】.  Directly computing \(P(o)\) is intractable, so agents minimize the **variational free energy** \(F\), which bounds surprise by Jensen’s inequality (i.e. \(F \ge -\ln P(o)\))【10†L43-L49】.  Formally, if \(Q(s)\) is an approximate posterior over hidden states \(s\), then
    \[
    F = \int Q(s)\ln\frac{Q(s)}{P(o,s)}\,ds
    = -\ln P(o) + \mathrm{KL}[Q(s)\,\|\,P(s|o)].
    \]
    At the optimal posterior (\(Q(s)=P(s|o)\)), \(F = -\ln P(o)\) (the surprise)【27†L470-L478】.  Thus minimizing \(F\) via belief updating or learning simultaneously maximizes model evidence.  Intuitively, **free energy** combines a fit term (log-evidence) with a complexity cost (Kullback-Leibler divergence)【27†L470-L478】, but crucially it depends only on observable data and internal beliefs, making it tractable to reduce【10†L43-L49】【27†L470-L478】.
    
    In a steady-state (homeostatic) regime, agents thus strive to minimize surprise (self-information) of sensory data.  This requires matching predictions to observations (accurate likelihoods and priors) or, if mismatches occur, rapidly updating beliefs.  The generative model (with its priors over states and likelihood mapping to observations) defines how surprise is evaluated: lower prior beliefs or narrower likelihoods both make an outcome less probable, raising surprise when that outcome happens【10†L115-L117】【51†L318-L324】.  We now explore in detail the factors that influence surprise in this framework.
    
    ## Factors Increasing Surprise
    
    - **Improbable or Unexpected Observations.**  Any observation \(o\) that the model assigns a low probability will produce high surprise \(I(o) = -\ln P(o)\)【10†L115-L117】.  For example, if a binary observation has prior probability 0.9 for one outcome and 0.1 for the other, observing the 0.1 outcome yields a large surprise.  In general, surprise grows without bound as \(P(o)\to 0\).  Thus *model mismatch* (when the true data-generating process differs from the agent’s model) causes unexpected outcomes with low predicted probability, boosting surprise.  This includes novel or outlier data that the model did not account for.
    
    - **Low Prior Probability.**  Even if an observation fits a likelihood, a very low prior belief on the hidden cause makes \(P(o)=\sum_s P(o|s)P(s)\) small, raising surprise.  In hierarchical models, higher-level priors influence lower-level predictions: a surprising event may occur simply because the prior expectation for its cause was low.  In effect, shrinking priors (more confident, narrow beliefs) amplify surprise for events lying in the tails of the prior.  Mathematically, since \(P(o)=\sum_s P(o|s)P(s)\), if \(P(s)\) is tiny for the true state, \(P(o)\) is also tiny, so \(-\ln P(o)\) is large.
    
    - **High Precision (Low Likelihood Variance).**  Precision is the inverse variance of a probability distribution.  A likelihood \(P(o|s)\) (e.g. Gaussian noise model) with small variance (high precision) is very “sharp”.  In such cases, even modest deviations between predicted and actual observations cause a large \((o-\text{prediction})^2/(2\sigma^2)\) term, making \(-\ln P(o|s)\) large.  For a Gaussian model \(o\sim N(\mu,\sigma^2)\), the surprise is
      \[
      I(o) = \frac{1}{2}\ln(2\pi\sigma^2) + \frac{(o-\mu)^2}{2\sigma^2}.
      \]
      With fixed prediction error \((o-\mu)\), decreasing \(\sigma\) (increasing precision) increases both the normalization term \(\ln\sigma\) and the squared-error term’s weight, raising surprise.  Thus **high confidence** in predictions (narrow likelihoods) makes violations more surprising【42†L406-L414】【51†L318-L324】.  Conversely, if the model assumes very high sensory noise (low precision), predictions are broad and any particular outcome is less surprising – although the baseline unpredictability (entropy) is higher.
    
    - **Mismatch in Noise Assumptions.**  When the agent’s assumed sensory noise (likelihood variance) is too low relative to reality, the true noisy observations will repeatedly lie far from the predicted mean. Each such outcome will have low \(P(o)\) under the naive model, leading to large prediction errors and surprise.  In contrast, if the agent correctly expects high noise, any given deviation is more expected (less surprising).  We will illustrate this with examples below.
    
    - **Prediction Errors and Precision-Weighting.**  In predictive coding formulations, each level of a hierarchical model generates top-down predictions that are compared with bottom-up signals, producing **prediction errors**.  Surprise accrues from these errors. Crucially, each prediction error is weighted by its precision (gain) before influencing higher levels.  High precision (attention) on a channel makes its errors more salient.  For instance, Feldman and Friston note that in a t-test the “prediction error” (difference of means) is divided by its standard error (precision)【42†L406-L414】, illustrating that a prediction error multiplied by precision drives surprise.  In neural terms, precision is encoded by synaptic gain【51†L318-L324】.  When precision weighting is high, even small discrepancies lead to large effective surprise.  Conversely, when precision is low (e.g. unattended inputs), errors are down-weighted and surprise is muted.  Hierarchically, expected precision at higher levels (e.g. belief about volatility) can modulate how much lower-level errors count as surprise.
    
    - **Action and Policy Effects.**  Agents can **act** to minimize expected future surprise.  Active inference casts action selection as minimizing **expected free energy** \(G(\pi)\) of policies \(\pi\).  \(G(\pi)\) includes terms for expected negative log-evidence (future surprise or “risk”) and expected information gain (epistemic value)【27†L490-L498】.  By choosing actions that lead to anticipated outcomes, agents reduce the surprise they eventually encounter.  When actions correctly fulfill predictions (e.g. moving to keep a preferred sensory state), surprise remains low.  **Action failure** – when an action produces an unpredicted outcome – yields a large immediate prediction error (surprise), prompting belief revision.  In essence, actions under active inference try to **realize** predictions (thus avoiding surprise); failure to do so forces updating of either predictions or future actions.
    
    - **Novelty and Exploration.**  Seeking novelty (for learning) can paradoxically increase surprise in the short term: exploring uncharted situations yields low prior \(P(o)\) and thus high surprise, but reduces future surprise by refining the model.  Conversely, “play it safe” actions that lead to well-known outcomes keep surprise low.  Some formulations (curiosity-driven models) explicitly balance surprise against information gain, but under strict free-energy principle agents prefer stimuli they can predict well.
    
    The table below summarizes these factors and their roles in the mathematics of surprise. (We use “+” to denote *increase* of surprise.)
    
    | **Factor**                              | **Mathematical Role**                                                                                                                                                                      | **Effect on Surprise**                                        |
    |-----------------------------------------|--------------------------------------------------------------------------------------------------------------------------------------------------------------------------------------------|---------------------------------------------------------------|
    | **Unlikely observation**                | Low predicted probability \(P(o)\) (e.g. outcome in tail of likelihood/prior) → large \(-\ln P(o)\)【10†L115-L117】.                                                                 | Large surprise (+).                                           |
    | **Model mismatch / bad likelihood**     | True generative process not captured by model → predicted \(P(o)\) misaligned with data, making actual \(P(o)\) small.                                                                     | High surprise (+).                                           |
    | **Low prior probability of cause**      | Small \(P(s)\) for true hidden state in the prior → small \(P(o)=\sum_sP(o|s)P(s)\).                                                                                                      | Higher surprise (+).                                         |
    | **Narrow likelihood (high precision)**  | Small variance \(\sigma^2\) in \(P(o|s)\) → normalization term \(\frac12\ln\sigma^2\) smaller but squared error term \((o-\mu)^2/(2\sigma^2)\) larger for mismatches.【42†L406-L414】  | Increases surprise if \(o\neq \mu\) (+).                     |
    | **Wide likelihood (low precision)**     | Large \(\sigma^2\) → broad prediction, small weight on any one error but larger baseline entropy.                                                                                         | Individual deviations less surprising (–), but *uncertainty* ↑ |
    | **High sensory noise (misestimated)**   | If model underestimates noise, true variability leads to frequent large \((o-\mu)\), raising surprise.                                                                                     | Surprise increases (+) if noise ignored.                     |
    | **Precision weighting (attention)**     | Prediction error \(\times\) precision (gain) → weighted surprise【42†L406-L414】【51†L318-L324】.                                                                                           | Amplifies surprise (+) for attended inputs.                  |
    | **Hierarchy (multi-level)**            | Higher-level priors shape lower-level expectations; unpredicted lower-level data propagate surprise upward.                                                                                 | Mismatch at any level adds to total surprise.               |
    | **Action (policy choice)**             | Policies choose actions to drive expected \(-\ln P(o)\) low (minimizing expected surprise)【27†L490-L498】. Failed actions → unexpected \(o\), spike in surprise.                          | Correct action reduces surprise; failed action causes spike. |
    | **Exploration / Novelty**             | Encountering new outcomes with initially low \(P(o)\).                                                                                                                                        | Increases surprise in short term (+), but allows learning.    |
    
    ## Likelihood, Prior, Posterior, and Precision
    Surprise depends on both **likelihoods** and **priors**.  The marginal likelihood \(P(o)\) is obtained by summing/integrating over hidden states \(s\): \(P(o)=\sum_s P(o|s)P(s)\).  Thus a low likelihood \(P(o|s)\) (for all \(s\)) or low prior \(P(s)\) (for the contributing \(s\)) makes \(P(o)\) small.  In Bayesian terms, if the *posterior* \(P(s|o)\) must place mass on unexpected states, this reflects high surprise in the data.  (Variational free energy can be written as \(F=-\ln P(o)+\mathrm{KL}[Q(s)\|P(s|o)]\); even if the agent’s chosen \(Q(s)\) is far from the true posterior, \(F\) remains an upper bound on surprise【10†L43-L49】【27†L470-L478】.)
    
    **Precision** (inverse variance) in both prior and likelihood critically modulates surprise.  For example, a strong prior (very narrow) punishes deviations more severely: if an outcome falls outside the narrow prior expectation, the prior factor \(P(s)\) is tiny for the true state and surprise skyrockets.  Similarly, high *likelihood* precision (low noise) makes any sensory prediction error count as a large surprise.  In predictive coding terms, the precision of a prediction error is the gain by which the brain weights that error【42†L406-L414】【51†L318-L324】.  When precision is high (focused attention), the same raw error leads to a larger effective surprise signal.  Crucially, precision itself can be context-sensitive and inferred (e.g., the brain may infer higher precision under stable conditions).  **Attention** is interpreted as inference about precision: attending to a modality raises its assumed precision【42†L406-L414】, so errors there become more surprising.
    
    By contrast, **high expected sensory noise** (low precision) means the agent expects variability; then particular mismatches are less surprising.  For a Gaussian likelihood, increasing the assumed noise variance \(\sigma^2\) spreads the distribution: the \((o-\mu)^2/(2\sigma^2)\) term shrinks, reducing surprise for a fixed error, at least up to a point.  In other words, when observations are inherently noisy, the agent lowers expectations of certainty, so large deviations do not carry as much negative log-probability.  Empirically, subjects with higher sensory uncertainty show broader tuning and smaller cortical responses to errors【51†L318-L324】. However, if the agent misestimates noise (e.g. assumes low noise but data are actually noisy), it will be repeatedly “surprised” by normal variability, causing erratic updates.
    
    **Bayesian surprise (information gain)** is a related concept: it measures the KL divergence between prior and posterior, quantifying how much beliefs must change.  While not identical to self-information, it reflects the *value* of surprise in learning.  High Bayesian surprise (large update) often accompanies high self-information: an unexpected outcome (low probability) forces a big belief revision.
    
    ## Hierarchical Models and Precision-Weighted Prediction Errors
    
    In a hierarchical generative model, each layer sends **predictions** downward and receives **prediction errors** from the layer below.  Surprise is then distributed across levels: an unexpected sensory input produces a large low-level error, which cascades upward as higher levels adjust.  Importantly, each prediction error at each level is weighted by its precision (which may depend on the inferred state at that level)【51†L318-L324】【53†L683-L690】.  This gives rise to the notion of *precision-weighted prediction errors* as the drivers of surprise.
    
    ```mermaid
    flowchart TD
        A[High-Level Beliefs (contexts)] --> B[Top-Down Prediction \n(of Lower States/Inputs)]
        B --> D[Prediction Error \n(Lower-Level)]
        C[Sensory Input (Data)] --> D
        D -->|weighted by precision| E[Prediction Error (Precision-Weighted)]
        E --> F[Belief Update at Higher Level]
    ```
    
    In this diagram, high-level beliefs generate predictions (B).  The actual sensory input (C) is compared with these predictions, yielding a raw prediction error (D).  This error is modulated by the precision (attentional gain) on the sensory channel, giving a weighted error (E).  That weighted error is used to update higher-level beliefs (F), thereby incorporating surprise into learning.  **Precision** (not shown explicitly) influences the arrow from D to E: higher precision means a steeper arrow (more weight).
    
    Because of this hierarchy, surprise can arise at any level.  For example, a low-level violation might be expected under current beliefs, producing moderate surprise.  However, if higher-level predictions were very confident (high precision), that same low-level deviation becomes a strong cue that the context belief is wrong, generating large higher-level surprise.  Conversely, if higher-level uncertainty is high (low precision), the system will be less “shocked” by low-level errors, attributing them to benign volatility.  Feldman and Friston illustrate this with a simple analogy: the statistical t-test is like a prediction error divided by an estimated standard error【42†L406-L414】.  The same prediction error can be either significant surprise (large t) or not, depending on its precision.
    
    Practically, this means surprise is **contextual**.  An odd sound in a quiet predictable environment is very surprising (high prediction error with high expected precision), but the same sound in a chaotic noisy crowd might not register as surprising (precision is low).  Recent experiments confirm that **precision-weighting can flip surprise**: for instance, unexpected interoceptive arousal (a shock or sudden heartbeat) can reduce the impact of sensory noise on perceived surprise or confidence【53†L683-L690】.
    
    ## Example Calculations of Surprise
    
    To ground these ideas, consider simple Gaussian examples.  Suppose an agent predicts a scalar outcome \(y\) with a Gaussian likelihood \(N(\mu,\sigma^2)\).  The surprise of observing \(y\) is:
    \[
    I(y) = -\ln\left(\frac{1}{\sqrt{2\pi\sigma^2}} e^{-\frac{(y-\mu)^2}{2\sigma^2}}\right)
    = \frac{(y-\mu)^2}{2\sigma^2} + \frac12\ln(2\pi\sigma^2).
    \]
    - **Changing prediction error**: Let \(\mu=0\), \(\sigma=1\).  If \(y=2\), then \(I=2.92\).  If \(y\) instead is 3 (larger error), \(I=4.05\).  Thus a larger deviation \((y-\mu)\) yields higher surprise.
    - **Changing precision**: Fix \(\mu=0\), \(y=2\), vary \(\sigma\).  For \(\sigma=2\), \(I\approx2.11\); for \(\sigma=1\), \(I\approx2.92\); for \(\sigma=0.5\), \(I\approx8.23\).  As \(\sigma\) decreases (precision increases), surprise rises steeply for the same error (see table).
    
    |  \( \mu=0, y=2\)   | \(\sigma=0.5\) | \(\sigma=1.0\) | \(\sigma=2.0\) |
    |--------------------|---------------|--------------|---------------|
    | \(I(y)\) (surprise) | 8.226         | 2.919        | 2.112         |
    
    This shows that the same outcome becomes much more surprising as the likelihood variance shrinks (precision grows).  Conversely, if the agent assumed \(\sigma\) was larger, surprise would be lower.
    
    - **Prior versus likelihood**: If instead the agent has a discrete prior \(P(s)\) over hidden states, an outcome \(y\) might only be possible under a rare state.  Suppose \(P(s_1)=0.01, P(s_2)=0.99\), and only \(s_1\) can produce a particular observation (so \(P(o|s_2)=0\)).  Then \(P(o)=0.01\) and \(I(o)=4.605\).  Raising the prior on \(s_1\) to 0.5 would reduce surprise to \(I(o)=0.693\).  Low priors on true causes sharply increase surprise.
    
    These calculations exemplify the roles in the table: small predicted probability (via low prior or narrow likelihood) produces large \(-\ln P(o)\).  Adjusting these parameters directly modulates surprise.
    
    ## Empirical and Simulation Evidence
    
    - **Sensory Noise and Surprise.**  Behavioral studies show that increasing sensory noise (variance) makes stimuli less predictable and typically reduces perceptual confidence【51†L318-L324】.  For example, Allen et al. (2016) found that adding noise to a visual motion task lowered subjects’ confidence【53†L683-L690】.  Intriguingly, they also showed that an unexpected arousal cue (a masked disgust face) *counteracted* this effect: arousal prediction errors reduced the impact of sensory noise on confidence, consistent with the view that interoceptive signals modulate precision on sensory inputs【53†L683-L690】.  In other words, when aroused, subjects treated the noisy stimulus as relatively more precise (less surprising) than without arousal.  This supports the role of precision weighting: context (arousal) changed the effective surprise of noisy data.
    
    - **Precision and Psychopathology.**  Computational studies link abnormal precision estimates to psychiatric symptoms.  Limongi et al. (2018) applied hierarchical drift-diffusion models to schizophrenia, finding that patients had **overly precise priors** relative to sensory evidence【35†L189-L193】.  In their task, patients prematurely acted on weak sensory evidence, consistent with over-weighting prior belief and under-weighting sensory input.  The dysconnection hypothesis predicts that impaired gain control (aberrant precision) leads to false inferences【35†L229-L237】.  Indeed, Limongi *et al.* conclude that schizophrenia involves “unduly precise prior beliefs to compensate for a failure of sensory attenuation”【35†L189-L193】.  This imbalance – too much confidence in internal predictions – causes normal sensory input to come as a “surprise,” contributing to hallucinations or delusions.
    
    - **Active Inference Simulations.**  Numerous modeling studies illustrate how surprise drives learning and action.  For example, Friston et al. (2015) showed that agents solving control tasks actively gather information to minimize expected free energy.  Reinforcement learning that accounts for surprise fits behavior better than pure reward maximization in some experiments【24†L1-L4】.  Although we omit detailed simulations here, these studies consistently demonstrate that agents reduce surprise by updating beliefs (learning the model) or by choosing actions that make sensory data more predictable.
    
    - **Perceptual Phenomena.**  The basic principles are evident in everyday perception.  Predictive coding models explain perceptual illusions (where priors dominate), repetition suppression (reduced response to predicted stimuli), and mismatch responses (surprise signals for unexpected inputs).  For instance, the mismatch negativity in EEG is larger for rare or unexpected tones.  These effects arise naturally from high surprise (prediction error) when an observation deviates from learned regularities.
    
    Overall, empirical findings support that **higher surprise corresponds to more cortical response, faster belief updating, and shifts in behavior**.  Precision manipulations (attention, pharmacology, arousal) systematically scale these surprise effects【53†L683-L690】【42†L406-L414】.
    
    ## Implications for Perception, Learning, and Psychopathology
    
    - **Perception.**  Surprise quantifies how much an input deviates from expectation.  In perception, minimizing surprise leads to efficient coding of regularities.  High surprise triggers enhanced processing of novel stimuli.  For example, an unexpected visual change (“oddball”) elicits strong neural responses, reflecting the large prediction error.  Attention (high precision) makes unexpected features more salient.  Conversely, in a predictable context, surprise is dampened and perception emphasizes predicted aspects.
    
    - **Learning.**  Learning (updating model parameters or structure) is driven by surprise.  Large prediction errors indicate model inadequacy, prompting substantial updates (e.g. large Bayesian learning rates).  Variational learning rules adjust synaptic weights to reduce future surprise.  In effect, the agent adapts priors and likelihoods to fit data, so that once learned, those outcomes become less surprising.  For example, encountering a consistent pattern of noisy inputs leads to increasing the assumed noise level, thereby reducing surprise over time.
    
    - **Psychopathology.**  Disorders like schizophrenia or autism have been linked to disrupted precision-scaling of prediction errors.  In schizophrenia, excessively high precision on priors (or low precision on sensory channels) means even mundane stimuli violate the brain’s predictions, yielding pathological surprise (hallucinations)【35†L229-L237】.  In autism, some models suggest reduced flexibility in precision, leading to either over- or under-reaction to sensory data.  Anxiety and depression may involve aberrant estimation of uncertainty, altering expected surprise.  Empirical studies (e.g. patient EEG/behavior) often find exaggerated or blunted surprise responses in these conditions, consistent with the active inference account.
    
    In summary, **surprise is a key driver of neural and behavioral change**.  By understanding the mathematical and computational factors that increase surprise, we gain insight into how agents perceive novelty, update beliefs, and sometimes malfunction when these mechanisms go awry.
    
    ## Comparison of Factors Affecting Surprise
    
    | **Factor**            | **Mathematical Role**                                             | **Qualitative Effect on Surprise**                              |
    |-----------------------|--------------------------------------------------------------------|----------------------------------------------------------------|
    | Low Prior on Cause    | Low \(P(s)\) for true state → low \(P(o)=\sum_sP(o|s)P(s)\)        | Increases surprise (outcome was unexpected)【10†L115-L117】.      |
    | Improbable Outcome    | Low \(P(o)\) directly, e.g. data far in likelihood tail           | Large surprise (self-information)【10†L115-L117】.               |
    | Narrow Likelihood     | Small variance \(\sigma^2\) in \(P(o|s)\)                         | Large \((o-\mu)^2/(2\sigma^2)\) term if \(o\neq\mu\)【42†L406-L414】. Increases surprise. |
    | Wide Likelihood       | Large \(\sigma^2\) (low precision)                                 | Lowers weight of deviations (reduces surprise), but raises baseline uncertainty. |
    | High Precision (Gain) | Error × precision in prediction error weighting                   | Amplifies surprise for a given error【42†L406-L414】【51†L318-L324】. |
    | High Noise (Low Precision)| Assumed \(\sigma^2\) large; predictions broad                  | Diminishes surprise for fixed error (because error is expected); increases entropy. |
    | Hierarchy (mismatch)  | Incorrect higher-level state → many lower-level errors            | Compound surprise across levels; large model update.           |
    | Action Success        | Outcome matches predicted reward/sensorium (high expected \(P(o)\))| Low surprise (goals achieved).                                  |
    | Action Failure        | Outcome deviates from expected; actual \(P(o)\) low              | High surprise; triggers re-planning.                           |
    
    ## How Surprise Arises and Is Resolved
    
    ```mermaid
    flowchart TD
        A[Agent with Generative Model] --> B[Make Prediction (based on priors)]
        B --> C[Receive Observation (Sensory Input)]
        C --> D{Predicted vs Actual}
        D -->|Mismatch| E[High Prediction Error (Surprise) ↑]
        D -->|Match| F[Low Prediction Error (Surprise) ↓]
        E --> G[Update Beliefs/Model or Select New Action]
        F --> G
        G --> H[Revised Model; Future Predictions Adjusted]
    ```
    
    This flowchart shows the *surprise timeline*.  The agent starts with a **prior model** (A) and makes **predictions** (B) about incoming data.  When an observation arrives (C), the agent compares it to its prediction.  If there is a **mismatch** (D), a high prediction error (surprise) is generated (E).  The agent then **updates its beliefs** or selects a corrective action (G) to reduce that error.  If the observation matches the prediction (D), surprise remains low (F), and little update is needed.  Over time, the updated model (H) will yield better predictions, resolving surprise.
    
    In summary, surprise arises whenever observations fall outside expected ranges; the agent responds by revising priors/likelihoods or altering actions.  Precision modulates each arrow: if the agent has high precision (confidence) at D, even a small mismatch leads to a large surprise node E.  If precision is low, D more likely follows the “Match” path.  The loop ensures that by acting or learning, future surprise is minimized.
    
    ## References
    - Definition of surprise: negative log-probability【10†L115-L117】. Variational free energy bounds surprise【10†L43-L49】【27†L470-L478】.
    - Precision concept: precision = inverse variance【51†L318-L324】; prediction errors are weighted by precision【42†L406-L414】【51†L318-L324】.
    - Aberrant precision in schizophrenia: high prior precision → false inferences【35†L189-L193】【35†L229-L237】.
    - Sensory noise and arousal (precision weighting): arousal prediction errors reduce noise-induced uncertainty【53†L683-L690】.
    - Feldman & Friston (2010) on attention as precision estimation【42†L406-L414】.
    - Parr & Friston (2017) on expected free energy and surprise【27†L470-L478】【27†L490-L498】.
    
 


    % \bibliographystyle{plain}
    % \bibliography{/home/donkarlo/Dropbox/repo/shared_research/bibliography/bibliography.bib}
\end{document}